%% oge-template.tex
%% Copyright 2026 CSharpKun & PositiveCHEEK
%
% This work may be distributed and/or modified under the
% conditions of the LaTeX Project Public License, either version 1.3
% of this license or (at your option) any later version.
% The latest version of this license is in
%   https://www.latex-project.org/lppl.txt
% and version 1.3c or later is part of all distributions of LaTeX
% version 2008 or later.
%
% This work has the LPPL maintenance status `maintained'.
% 
% The Current Maintainer of this work is CSharpKun

\documentclass{exam}

\usepackage{fontspec}
\usepackage{polyglossia}

\setdefaultlanguage{russian}
\setmainfont{Liberation Serif} 

\author{CSharpKun \& PositiveCHEEK}
\date{2026}

\makeatletter
	\footer{}{© {\@date} \@author}{}
\makeatother

\newcommand{\task}[1]{%
	\noindent
	\makebox[0pt][r]{%
		\fbox{\makebox[0.8cm][c]{\textbf{#1}}}%
		\hspace{0.2cm}%
	}%
}

\newcommand{\descbox}[1]{
	\fbox{
		\parbox{1.0\textwidth}{ 
			#1
		}
	}
}

\renewcommand{\thechoice}{\arabic{choice}}

\setlength{\parindent}{25pt}

\begin{document} % НАЧАЛО ДОКУМЕНТА

\begin{center}
	\textbf{\large{Основной государственный экзамен по [ПРЕДМЕТ]}}
\end{center}

\begin{center}
	\textbf{\large Инструкция по выполнению работы}
\end{center}


Экзаменационная работа состоит из двух частей, включающих в себя
[КОЛИЧЕСТВО ЗАДАНИЙ] заданий. Часть 1 содержит [КОЛИЧЕСТВО ЗАДАНИЙ 1 ЧАСТИ] заданий, часть 2 содержит [КОЛИЧЕСТВО ЗАДАНИЙ 2 ЧАСТИ] заданий.

На выполнение экзаменационной работы по музыке отводится
[ВСЕГО ЧАСОВ] ([ВСЕГО МИНУТ]).

% Напишите ваши требования к ответам на задачи

Ответами к заданиям части 1 является число или
последовательность цифр. Ответ запишите в поле ответа в тексте работы,
а затем перенесите в бланк ответов № 1.

Решения заданий части 2 и ответы к ним запишите на бланке
ответов № 2. Задания можно выполнять в любом порядке. Текст задания
переписывать не надо, необходимо только указать его номер.

Все бланки заполняются яркими чёрными чернилами. Допускается
использование гелиевой или капиллярной ручки.

Сначала выполняйте задания части 1. Начать советуем с тех заданий,
которые вызывают у Вас меньше затруднений, затем переходите к другим
заданиям. Для экономии времени пропускайте задание, которое не удаётся
выполнить сразу, и переходите к следующему. Если у Вас останется время,
Вы сможете вернуться к пропущенным заданиям.

При выполнении части 1 все необходимые вычисления,
преобразования выполняйте в черновике. \textbf{Записи в черновике, а также
в тексте контрольных измерительных материалов не учитываются при
оценивании работы.}

% Убрать если рисунков нет

Если задание содержит рисунок, то на нём непосредственно в тексте
работы можно выполнять необходимые Вам построения. Рекомендуем
внимательно читать условие и проводить проверку полученного ответа.

% Убрать если справочных материалов нет

При выполнении работы Вы можете воспользоваться справочными
материалами, выданными вместе с вариантом КИМ.

Баллы, полученные Вами за выполненные задания, суммируются.
Постарайтесь выполнить как можно больше заданий и набрать наибольшее
количество баллов.

После завершения работы проверьте, чтобы ответ на каждое задание
в бланках ответов № 1 и № 2 был записан под правильным номером.

\medskip

\begin{center}
	\textit{\textbf{Желаем успеха!}} 
\end{center}

\newpage

\begin{center}
	\textbf{\large{Часть 1}}
\end{center}

% Изменить под ваши требования к первой части

\begin{center}
\descbox{
	\textbf{
		\textit{
			Ответами к заданиям 1–[ВСЕГО ЗАДАНИЙ] являются число или последовательность
			цифр, которые следует записать в БЛАНК ОТВЕТОВ № 1 справа от
			номера соответствующего задания, начиная с первой клеточки. Если
			ответом является последовательность цифр, то запишите её \underline{без
			пробелов и других дополнительных символов}. Каждый символ пишите
			в отдельной клеточке в соответствии с приведёнными в бланке
			образцами.
		}
	}
}
\end{center}

\task{1}
[ТЕКСТ ПЕРВОГО ЗАДАНИЯ]
\bigskip
\bigskip

\task{2}
[ТЕКСТ ВТОРОГО ЗАДАНИЯ]

\begin{oneparchoices}
	\choice [ВЫБОР]
	\choice [ВЫБОР]
	\CorrectChoice [ПРАВИЛЬНЫЙ ВЫБОР]
	\choice [ВЫБОР]
\end{oneparchoices}

\bigskip
\bigskip

\task{3}
[ТЕКСТ ТРЕТЬЕГО ЗАДАНИЯ]

\end{document}
